% Chapter 2: Finite differences from Taylor's theorem (fundamentals author). PLAN.md section 1.4.
% Uses F2 (fig:fd-slopes), S2 (sb:remainder), E2 (code/02-taylor-steps.cpp), the margin table tab:two-point.
% Language pass for students (PLAN.md section 12.9): the wording only; content, evidence and order as at 3192e8a.
\chapter[Finite differences from Taylor's theorem]{Finite differences from\\ Taylor's theorem}\label{ch:taylor}

\Cref{fig:secant} showed that the secant through $x - h$ and $x + h$ lies closer to the tangent than the secant
through $x$ and $x + h$. This chapter works out by how much. We need one tool for it: Taylor's theorem, which you
know from calculus. With it we derive the two simplest finite differences and the error that each one makes. Then we
find a rule that predicts how that error changes when we change the step. All of it can be done with pencil and
paper, and we check every prediction by measuring it on the heat capacity of nitrogen.

\section{The secant and the tangent}\label{sec:taylor-secant}

We start with the secant through $x$ and $x + h$. Its slope is the \emph{forward difference}
\begin{equation}
  f'(x) \approx \frac{f(x + h) - f(x)}{h}.
  \label{eq:forward}
\end{equation}
How far is it from $f'(x)$? Taylor's theorem answers that. If $f$ has a continuous second derivative between $x$
and $x + h$, then
\begin{equation*}
  f(x + h) = f(x) + h f'(x) + \frac{h^2}{2} f''(\xi)
\end{equation*}
for some point $\xi$ between $x$ and $x + h$.\cite{burden2016} The first two terms on the right are the tangent line.
The last term is the \emph{remainder}: the price of stopping the series after the term in $h$. We do not know $\xi$.
We only know that it lies somewhere between $x$ and $x + h$. Now subtract $f(x)$ from both sides and divide by $h$:
\begin{equation}
  \frac{f(x + h) - f(x)}{h} = f'(x) + \frac{h}{2} f''(\xi).
  \label{eq:forward-error}
\end{equation}
The left-hand side is the forward difference. So the forward difference is $f'(x)$ plus an error of
$h f''(\xi) / 2$. This error is the \emph{truncation error}. It is the error of the formula itself, even when every
number is computed exactly. The name says where it comes from: we truncated the Taylor series. The picture tells the
same story. A secant differs from the tangent because the curve bends, and the bending is $f''$.

Let us apply this to the enthalpy of nitrogen. There $\enth'' = \dd\cp / \dd T$, and at $T_0$ the closed form gives
\SI{\valCpSecondExact}{\joule\per\mole\per\kelvin\squared}. This is positive: the heat capacity still rises at
\SI{1000}{\kelvin}. So the forward difference comes out too high, by about half of $\enth''$ times the step.

\begin{margintable}
\caption{The two rules of this chapter, as $f'(x) \approx \sum_k w_k f(x + n_k h) / h$. The complete tables are in
\cref{ch:tables} (\cref{tab:weights-central-first}).\label{tab:two-point}}
\centering\footnotesize\setlength{\tabcolsep}{4pt}
\begin{tabular}{@{}lll@{}}
\toprule
 & \code{Forward<1>} & \code{Central<2>}\\
\midrule
nodes $n_k$ & $0, 1$ & $-1, 1$\\
weights $w_k$ & $-1, 1$ & $-\frac{1}{2}, \frac{1}{2}$\\
error & $\frac{h}{2} f''$ & $\frac{h^2}{6} f'''$\\
\bottomrule
\end{tabular}
\end{margintable}

The library calls this rule \code{Forward<1>}.\mnote{\code{P} in \code{Forward<P>} and \code{Central<P>} is the
accuracy order, not the number of points: \code{Central<2>} uses two points for $f'$, \code{Forward<1>} two.} The
number in the name is the rule's accuracy order, which \cref{sec:taylor-order} defines. The library, and this book,
write a rule as two lists. The \emph{nodes} $n_k$ are the offsets of the sample points, in units of $h$. The
\emph{weights} $w_k$ say how much each sample counts. The estimate is then $\sum_k w_k f(x + n_k h) / h$. For the
forward difference the nodes are 0 and 1 and the weights $-1$ and 1, which gives back $(f(x + h) - f(x)) / h$.
\Cref{tab:two-point} gives the nodes and weights of the two rules of this chapter.

\section{Centring cancels a term}\label{sec:taylor-central}

The forward difference samples one side of $x$ only. Now we use both sides. We expand $f(x + h)$ and $f(x - h)$, this
time one term further. If $f$ is smooth enough, that is, if it has enough continuous derivatives,
\begin{align*}
  f(x + h) &= f(x) + h f'(x) + \frac{h^2}{2} f''(x) + \frac{h^3}{6} f'''(x) + \frac{h^4}{24} f^{(4)}(x) + \cdots\\
  f(x - h) &= f(x) - h f'(x) + \frac{h^2}{2} f''(x) - \frac{h^3}{6} f'''(x) + \frac{h^4}{24} f^{(4)}(x) - \cdots
\end{align*}
Compare the two lines term by term. The terms with even powers of $h$ are the same in both lines. The terms with odd
powers differ only in sign. So when we subtract the second line from the first, $f(x)$, $f''$ and $f^{(4)}$ drop out,
and the remaining terms double. Dividing by $2h$ gives
\begin{equation}
  \frac{f(x + h) - f(x - h)}{2h} = f'(x) + \frac{h^2}{6} f'''(x) + \frac{h^4}{120} f^{(5)}(x) + \cdots
  \label{eq:central}
\end{equation}
The left-hand side is the \emph{central difference}. It is the library's \code{Central<2>}, and its default rule.
Compare its error with that of the forward difference: the term with $h$ has gone. The bending $f''$, which spoils
the forward difference, affects the two sides equally and cancels. What is left starts with $h^2$ and involves
$f'''$, the rate at which the bending changes. We can also write the error with a remainder at an unknown point
$\xi$, in Lagrange's form, as we did for the forward difference, instead of as an infinite series. The result reads
\begin{equation}
  \frac{f(x + h) - f(x - h)}{2h} = f'(x) + \frac{h^2}{6} f'''(\xi), \qquad x - h < \xi < x + h,
  \label{eq:central-error}
\end{equation}
which holds whenever $f'''$ is continuous on the interval. The sidebar below proves it, and shows why the next term
is of order $h^4$ (\cref{sb:remainder}).

The factor $h^2 / 6$ is positive, so the sign of the error is the sign of $f'''$. The central difference
overestimates $f'$ where $f''' > 0$, and underestimates it where $f''' < 0$. For nitrogen at \SI{1000}{\kelvin},
$\enth''' = \dd^2 \cp / \dd T^2$ is negative: the heat capacity rises, but ever more slowly. So the central difference
comes out slightly low.

Note what the central difference does not use: the value $f(x)$, which has weight zero (\cref{tab:two-point}). So it
costs two evaluations of $f$, the same as the forward difference, and it is far more accurate for small $h$. That is
why it is the default.

There is a second way to see the cancellation. The mirror image of the forward difference is the \emph{backward
difference} $(f(x) - f(x - h)) / h$. The library calls it \code{Backward<1>}; it is the rule for a function known
only below $x$. The same expansion gives its error as $-h f''(x) / 2 + \order(h^2)$.\mnote{$\order(h^2)$, ``of order
$h^2$'': terms of size $h^2$ and smaller, which shrink at least as fast as $h^2$ (\cref{sec:taylor-order}).} This is
the same size as the error of the forward difference, with the opposite sign. The central difference is exactly the
mean of the two,
\begin{equation*}
  \frac{f(x + h) - f(x - h)}{2h}
  = \frac{1}{2} \left( \frac{f(x + h) - f(x)}{h} + \frac{f(x) - f(x - h)}{h} \right),
\end{equation*}
so their errors of order $h$ cancel, and what remains is of order $h^2$.

We now have two rules and know their errors: an error of order $h$ for the forward difference, and of order $h^2$ for
the central difference. The next section turns these powers of $h$ into a prediction that we can test.

\begin{sidebar}[label=sb:remainder]{Where the constant 1/6 comes from}
Taylor's theorem with the remainder in Lagrange's form says the following. If $f$ has $n + 1$ continuous derivatives
between $x$ and $x + h$, then
\begin{equation*}
  f(x + h) = \sum_{j=0}^{n} \frac{h^j}{j!} f^{(j)}(x) + \frac{h^{n+1}}{(n+1)!} f^{(n+1)}(\xi)
\end{equation*}
for some $\xi$ strictly between $x$ and $x + h$. The same holds with $-h$ in place of $h$.\cite{burden2016} We take
$n = 2$, once on each side of $x$:
\begin{equation*}
  f(x \pm h) = f(x) \pm h f'(x) + \frac{h^2}{2} f''(x) \pm \frac{h^3}{6} f'''(\xi_\pm),
\end{equation*}
with $\xi_+$ in $(x, x + h)$ and $\xi_-$ in $(x - h, x)$. We subtract the two and divide by $2h$:
\begin{equation*}
  \frac{f(x + h) - f(x - h)}{2h} = f'(x) + \frac{h^2}{6} \cdot \frac{f'''(\xi_+) + f'''(\xi_-)}{2}.
\end{equation*}
The mean of $f'''(\xi_+)$ and $f'''(\xi_-)$ lies between the two values. A continuous function takes every value
between two of its values; this is the intermediate value theorem. So the mean equals $f'''(\xi)$ for some $\xi$ in
$(x - h, x + h)$, and that is \cref{eq:central-error}. The $1/6$ is the $1/3!$ of the cubic term. It survives
unchanged because the two cubic terms add: their sum carries a factor 2, which the division by $2h$ removes.

Why is the next term of order $h^4$? In the series for $f(x + h) - f(x - h)$, every even power of $h$ cancels. Only
the odd powers $h, h^3, h^5, \ldots$ remain. After the division by $2h$, only the even powers $h^0, h^2, h^4,
\ldots$ are left. So the error of the central difference is a series in $h^2$. The forward difference has no such
symmetry, and its error contains every power of $h$.
\end{sidebar}

\section{Order of accuracy}\label{sec:taylor-order}

This section gives the power of $h$ in an error term a name, the accuracy order, and shows what it predicts. Both
error terms have the same form: a constant, times a power of $h$, times a derivative of $f$. For the forward
difference this is $\frac{1}{2} h^1 f''$, and for the central one $\frac{1}{6} h^2 f'''$. In general a rule for the
$N$-th derivative has a truncation error whose leading term, the one that dominates for small steps, is
\begin{equation*}
  \pm \errc\, h^{\aorder} f^{(\aorder + \dorder)}(x),
\end{equation*}
where $\errc > 0$ is the size of the term. Its sign is fixed by the rule: $+$ for both rules here, and
\cref{ch:tables} gives it for every rule. The power $\aorder$ is the rule's \emph{accuracy order}.\mnote{The
\emph{derivative order}
$\dorder$ says which derivative is computed; this chapter has $\dorder = 1$ throughout. \Cref{ch:orders} treats
$\dorder = 2$ to 4.} The forward difference has accuracy order 1, and the central difference accuracy order 2. We
write that their errors are $\order(h)$ and $\order(h^2)$.

This is a statement about small steps. As $h$ tends to 0, the error divided by $h^{\aorder}$ stays bounded, and the
leading term describes the error better and better. For larger steps the next terms matter too. For the central
difference the next term is of order $h^{\aorder + 2}$, here $h^4$, as \cref{eq:central} shows.

In practice the accuracy order gives us a prediction. Suppose the leading term dominates. When we halve the step,
$h^{\aorder}$ becomes $h^{\aorder} / 2^{\aorder}$, so the error is divided by $2^{\aorder}$. Halving the step
therefore halves the error of the forward difference, and quarters that of the central difference. To gain a factor
of 10 in accuracy, the forward difference needs a step 10 times smaller. The central difference needs a step only
$\sqrt{10} \approx 3.2$ times smaller.

\begin{figure}[t]
\sidecaption{The relative error of \code{Forward<1>} and \code{Central<2>} against the step $h$, for the heat
capacity of nitrogen, $\dd\enth / \dd T$ at $T_0 = \SI{1000}{\kelvin}$, at large steps. On logarithmic axes the errors
fall on lines of slope 1 and 2: halving the step halves one error and quarters the other. Below \SI{0.01}{\kelvin}
the central difference stops improving; \cref{ch:rounding} explains why.\label{fig:fd-slopes}}
\centering
% F2 fig:fd-slopes (Chapter 2, text width): relative error of the forward difference (Forward<1>) and the central
% difference (Central<2>) against the step on the thread, H(T) at T0 = 1000 K, for the large steps 1e2 to 1e-2 K
% (DS1, data/step-error-cp.csv). On log axes the errors fall on lines of slope 1 and 2 while truncation dominates.
% The range stops at 1e-2 K, where the central difference stops improving (rounding, Chapter 3): the caption points
% ahead (editor, stage 2, decision S2-14). A domain restriction on a log axis is in natural logarithms:
% ln(1e-2) = -4.61, ln(1e2) = 4.61.
\begin{tikzpicture}
\begin{axis}[eqtufte, width=95mm, height=45mm, xmode=log, ymode=log,
    xlabel={step $h$ (K)}, ylabel={relative error},
    xmin=1e-2, xmax=1e2, ymin=1e-12, ymax=1e-1,
    xtickten={-2,-1,0,1,2}, ytickten={-12,-10,-8,-6,-4,-2},
    restrict x to domain=-4.62:4.62]
  \addplot[muted, mark=*] table[col sep=comma, x=h, y=forward1, comment chars={\#}] {data/step-error-cp.csv};
  \addplot[accent, mark=*] table[col sep=comma, x=h, y=central2, comment chars={\#}] {data/step-error-cp.csv};
  \node[muted, font=\sffamily\scriptsize, anchor=south east, align=right] at (axis cs:3, 1e-3)
    {forward difference: slope 1,\\ halving $h$ halves the error};
  \node[accent, font=\sffamily\scriptsize, anchor=north west, align=left] at (axis cs:0.3, 1e-9)
    {central difference: slope 2,\\ halving $h$ quarters the error};
\end{axis}
\end{tikzpicture}

\end{figure}

We can check the prediction by hand. Take $f(x) = x^3$ at $x = 1$, where $f'(1) = 3$ and $f'' = f''' = 6$. We use
the steps $0.1$ and $0.05$:
\begin{center}\small
\begin{tabular}{@{}lcccc@{}}
\toprule
 & $f(1 + h)$ & $f(1 - h)$ & forward, error & central, error\\
\midrule
\tablebody{hand-cube}
\bottomrule
\end{tabular}
\end{center}
Halving the step divides the forward error by a little more than 2, and the central error by exactly 4. The error
terms predict the sizes too. For the forward difference, $h f''/2 = 3h$ gives $0.3$ and $0.15$; the small rest is
the next term, $h^2$. For the central difference, $h^2 f'''/6 = h^2$ gives $1/100$ and $1/400$. Here the prediction
is exact, because the fourth and higher derivatives of a cubic vanish.

How small must $h$ be for the leading term to dominate? It must be small compared with the distance over which $f$
changes its character. Then the higher derivatives in the next terms do not outweigh the extra powers of $h$ that
they come with. The enthalpy of nitrogen is a smooth function of $T/\SI{1000}{\kelvin}$, so steps of many kelvin are
already small in this sense. Most engineering functions vary on a scale comparable to the size of their argument.
\Cref{ch:rounding} builds the library's default step on that observation, and shows what to do when it fails.

\Cref{fig:fd-slopes} tests the prediction on the nitrogen example, $\enth(T)$ at $T_0 = \SI{1000}{\kelvin}$, with the
closed form of $\cp$ as the answer key. Both axes of the figure are logarithmic. Taking logarithms turns an error
$\errc h^{\aorder}$ into $\log \errc + \aorder \log h$, so on such axes it is a straight line of slope $\aorder$. The
two measured curves are such lines, of slopes 1 and 2. At every step shown, the central difference is more accurate
by orders of magnitude.

\needspace{22\baselineskip}% the lead-in and the whole listing on one page, not split around Figure 2.1
The same check takes a few lines of code. The listing computes both differences at steps of \SI{2}{\kelvin} and
\SI{1}{\kelvin}. For each rule it prints the errors at the two steps and their ratio:
\snippet{02-taylor-steps}{forward-vs-central}
The forward ratio comes out close to 2, and the central ratio close to 4, as predicted. At such steps the rounding
errors of \cref{ch:rounding} are negligible, and the next terms of the series shift the ratios only slightly.

\section{Exactness on polynomials}\label{sec:taylor-exact}

The error terms give a second check, an exact one that needs no computer. The forward difference has the error
$h f''(\xi) / 2$, which vanishes when $f'' = 0$. So the forward difference is exact for every straight line. It must
be, since a secant of a line is the line itself. The central difference has the error $h^2 f'''(\xi) / 6$, which
vanishes when $f''' = 0$. So it is exact for every quadratic. We can confirm this with pencil and paper. For
$f(x) = a + b x + c x^2$,
\begin{equation*}
  \frac{f(x + h) - f(x - h)}{2h} = \frac{2bh + 4cxh}{2h} = b + 2cx = f'(x)
\end{equation*}
for every $h$. For a cubic the central difference is not exact. For $f(x) = x^3$ it gives $3x^2 + h^2$, and the
excess $h^2$ is the error term $h^2 f''' / 6$ with $f''' = 6$. In general, a rule of accuracy order $\aorder$ for the
$\dorder$-th derivative is exact for polynomials of degree up to $\aorder + \dorder - 1$ (\cref{ch:stencils}). For
the central difference, $\aorder = 2$ and $\dorder = 1$ give degree 2: the quadratics we just checked. ``Exact'' here
means exact in exact arithmetic. A computer still rounds the values, and that is the subject of the next chapter.

The structure of \cref{eq:central} also shows where more accuracy can come from. Write $\Dc{h}$ for the central
difference at step $h$. Its error is $\frac{1}{6} h^2 f''' + \order(h^4)$. Since $(h/2)^2 = h^2 / 4$, the error of
$\Dc{h/2}$ has a quarter of that $h^2$ term. So $4 \Dc{h/2}$ has the same $h^2$ term as $\Dc{h}$, and the combination
\begin{equation*}
  \frac{4 \Dc{h/2} - \Dc{h}}{3} = f'(x) + \order(h^4)
\end{equation*}
cancels the $h^2$ term: the subtraction leaves $3 f'(x)$ plus terms of order $h^4$, and the division by 3 leaves
$f'(x)$ with an error of order $h^4$. This combination is the four-point rule \code{Central<4>} of \cref{ch:stencils} in disguise
(\cref{sb:richardson-disguise}). Applied again and again, the same idea is the heart of Ridders' method
(\cref{ch:ridders}).

\section{What this chapter did not say}\label{sec:taylor-not-said}

This chapter derived two rules, their errors and the accuracy order, all from Taylor's theorem. Every formula in it
rests on three assumptions. Here they are, with the places where the book takes each one up.
\begin{itemize}
  \item \emph{Exact arithmetic.} In a computer the values of $f$ are rounded. \Cref{fig:fd-slopes} stops at
    \SI{0.01}{\kelvin} because below that step the central difference stops improving. Why, and how to choose the
    step, is the subject of \cref{ch:rounding}.
  \item \emph{Two points.} More points buy a higher accuracy order, and one-sided rules of higher order exist
    (\cref{ch:stencils}). Higher derivatives follow in \cref{ch:orders}.
  \item \emph{Smoothness.} The error terms need $f''$ or $f'''$ to exist and to be continuous on the whole interval
    that the rule samples. If $f$ has a kink or a jump inside that interval, the error terms say nothing. The result
    can then be wrong without any warning (\cref{sec:using-pitfalls}).
\end{itemize}
